\chapter{\'Etica y protecci\'on de datos}
\label{cap:etica}

Este cap\'itulo resume \texttt{docs/etica/ETHICS.md}, documento de
referencia completo para cualquier aspecto \'etico o de gesti\'on de datos
de BIOPET no cubierto aqu\'i.

\section{Uso de datos acad\'emicos}

BIOPET usa exclusivamente dos categor\'ias de datos, ambas generadas por el
propio equipo: datos sint\'eticos de desarrollo/prueba (usuario
administrador sembrado, datos de registro creados durante las pruebas) y
datos emp\'iricos de evaluaci\'on generados por herramientas automatizadas
(k6, JaCoCo, evidencia OWASP) contra el propio sistema, sin intervenci\'on
de personas ni datos personales de terceros.

\section{Ausencia de datos cl\'inicos reales}

El sistema, en el alcance implementado hasta \texttt{v0.9.0-rc}, no
almacena historiales m\'edicos ni datos cl\'inicos de mascotas: \'unicamente
nombre, especie, raza, fecha de nacimiento y el identificador del due\~no.
No se ha desplegado con datos de clientes reales de una cl\'inica
veterinaria a la fecha de este informe.

\section{Credenciales de desarrollo}

La cuenta administradora sembrada autom\'aticamente
(\texttt{admin@biopet.ec}) y las credenciales de ejemplo del README son
datos ficticios de desarrollo, documentados como tales en el propio
repositorio, y no corresponden a ninguna persona real. Este informe
\textbf{no reproduce el valor literal de ninguna contrase\~na}, siguiendo
la misma pr\'actica ya aplicada en \texttt{docs/mediciones/sec/} y
\texttt{docs/postman/}.

\section{No inclusi\'on de secretos}

Ning\'un documento de este informe, ni sus figuras (cuando se agreguen),
debe contener: contrase\~nas con valor, JSON Web Tokens completos
(\texttt{eyJ...}), valores de cookies, el encabezado
\texttt{Authorization: Bearer} con valor real, claves privadas ni el
identificador \texttt{jti} de un token real. Esta regla se valid\'o de
forma expl\'icita al cierre de esta fase (secci\'on de validaciones,
cap\'itulo final de este informe).

\section{Protecci\'on de JWT y cookies}

El dise\~no del sistema (cap\'itulo~\ref{cap:resultados-jaime}) evita por
construcci\'on la exposici\'on de JWT: cookies \texttt{HttpOnly} (no
legibles desde JavaScript del navegador), \texttt{Secure} (solo sobre
HTTPS) y \texttt{SameSite=Strict}. La documentaci\'on de este informe sigue
la misma disciplina: cuando se describe un atributo de cookie, nunca se
reproduce su valor.

\section{Logs sin contrase\~nas}

\texttt{AuthenticationAuditService} (cap\'itulo~\ref{cap:resultados-jaime})
solo acepta \texttt{ip} y \texttt{subject} como par\'ametros en cada
m\'etodo p\'ublico: estructuralmente, no existe ninguna v\'ia de c\'odigo
para que una contrase\~na, un JWT completo o un encabezado
\texttt{Authorization} lleguen al logger.

\section{Responsabilidad compartida y l\'imites del entorno acad\'emico}

Los tres integrantes del equipo comparten la responsabilidad de no
introducir datos personales identificables ni secretos en el repositorio
p\'ublico. El sistema se declara expl\'icitamente como un artefacto
acad\'emico: si en una fase posterior se procesaran datos reales de
clientes de una cl\'inica veterinaria, \texttt{docs/etica/ETHICS.md}
deber\'ia actualizarse para incorporar una base legal de tratamiento y,
si aplica en Ecuador, referencia a la Ley Org\'anica de Protecci\'on de
Datos Personales (LOPDP) --- actualizaci\'on que no corresponde a esta
entrega por no existir todav\'ia ese procesamiento.

\section{Consentimiento informado para la prueba SUS}

Cuando se ejecute la prueba de usabilidad (SUS, m\'inimo diez
participantes externos al equipo), cada participante debe firmar el
formulario de consentimiento informado disponible en
\texttt{docs/etica/plantilla.md} antes de realizar la tarea de referencia.
Los formularios firmados no se suben al repositorio p\'ublico; los
resultados agregados se identifican \'unicamente por c\'odigo de
participante (P01, P02, \ldots).

Para el detalle completo de cada punto, ver \texttt{docs/etica/ETHICS.md}.
